
\svnid{$Id: $}
\svnid{$URL: $}

\section{REVIEW MISSING: Curved channel reference for junction advection testcases}
\newrefsegment

%---------------------------------------------------------------------------------
\section*{Quality Assurance}
\begin{tabular}{@{}p{27mm}@{}p{0mm}p{\textwidth-52mm-48pt}p{15mm}p{10mm}}
\textbf{Date} & \textbf{:} & May 22 2019  \\ 
\textbf{Author} & \textbf{:} & Rinske Hutten  & \textbf{Initials:} & RH \\
\textbf{Review} & \textbf{:} &   & \textbf{Initials:} & \ldots 
\end{tabular}

%---------------------------------------------------------------------------------
\section*{Version information}
{{
\begin{tabular}{@{}p{27mm}@{}p{0mm}p{\textwidth-27mm-24pt}}
\textbf{Executable}   & \textbf{:} & Deltares, D-Flow FM Version 1.2.54.64224M, Jun 28 2019, 15:14:31  \\
\textbf{Location input}     & \textbf{:} & \url{https://repos.deltares.nl/repos/DSCTestbench/trunk/cases/e02_dflowfm/f104_1D_numerical-aspects/c16_curved-channel-junction-advection-reference} \\
\textbf{Revision input} & \textbf{:} & - \\
\textbf{Location output}  & \textbf{:} & \url{https://repos.deltares.nl/repos/DSCTestbench/trunk/references/win64/e02_dflowfm/f104_1D_numerical-aspects/c16_curved-channel-junction-advection-reference/DFM_OUTPUT_Flow1D} \\
\textbf{Revision output} & \textbf{:} & -
\end{tabular}
}}

%---------------------------------------------------------------------------------
\section*{Purpose}
The purpose of the test \textbf{e02}(dflowfm)\_\textbf{f104}(numerical-aspects)\_\textbf{c16}(curved-channel-junction-advection-reference) is to be used as a reference case for junction advection test cases by showing that the differences of observed waterdepths and water levels at the connection node and the velocities near the connection nodes are due to the directional change of the channels. 

%---------------------------------------------------------------------------------
\section*{Linked claims}
Channels with a directional change at the connection node give higher water depths and water levels at the connection node for steady state conditions (flow diversion and acceleration at junction; equidistant grid).


%---------------------------------------------------------------------------------
\section*{Approach}
Water depths and velocities of two systems with a directional change are compared with the results of two systems with straight channels. More precisely, results of sub-models \textbf{\textcolor[rgb]{1,0,0}{T3}} and \textbf{\textcolor[rgb]{1,0,0}{T4}} are validated against the results of sub-model \textbf{\textcolor[rgb]{1,0,0}{T1}} and \textbf{\textcolor[rgb]{1,0,0}{T2}}(see \autoref{fig:e02_f104_c16_layout}).  
\begin{figure}[H]
    \centering
    \includegraphics*[width=1.0\textwidth]{figures/e02_f104_c16_layout.jpg}
    \caption{Lay-out of the test model}
    \label{fig:e02_f104_c16_layout}
\end{figure} 

%---------------------------------------------------------------------------------
\section*{Conclusion} 
A comparison of the results of this reference case resulted in higher water depths and water levels at the connection nodes for curved channels with steady state conditions (flow diversion and acceleration at junction; equidistand grid). A smaller angle (closer to 90\textdegree \,  instead of 180\textdegree) between two branches gave even higher water depths and water levels. Next to this, the curved channels had higher velocities before the connection nodes and lower velocities behind the connection nodes compared to straight channels. The magnitude of these differences between straight and curved channels depends in similar way as the water depths on the angle between two branches. The discharge was the same for all systems.
 
%---------------------------------------------------------------------------------
\section*{Model description}
he test comprises of four separate sub-models (see \autoref{fig:e02_f104_c16_layout}), all having an \underline{equidistant} grid-spacing of 1000 m.
\begin{enumerate}
	\item Sub-model \textbf{\textcolor[rgb]{1,0,0}{T1}} comprises of one branch \textbf{B1} (length 25000 m) only (see \autoref{fig:e02_f104_c16_layout}).
		\begin{itemize}
			\item At T1\_Bnd\_B1\_x0m a constant discharge of 790.57 $m^3$/s is imposed. 
			\item At T1\_Bnd\_B1\_x25000m a constant water level of 1.0 m is imposed.
			\item Branch B1 contains three "open Y-Z rectangular cross-sections" (width=\textcolor[rgb]{0.071,0,1}{\textcolor[rgb]{0.5,0,0.15}{\textbf{500}}} m, Chezy=50 $m^{0.5}$/s), respectively having an elevation of \textbf{10.15} m AD at T1\_Bnd\_B1\_x0m, \textbf{10.00} m AD at fixed water-level-point T1\_h\_pnt\_x15000m, and \textbf{0.00} m AD at \newline T1\_Bnd\_B1\_x25000m.
		\end{itemize} \


	\item Sub-model \textbf{\textcolor[rgb]{1,0,0}{T2}} comprises of two branches \textbf{B1} (length 15000 m) and \textbf{B2} (length 10000 m), that are mutually connected at connection node T2\_ConNode (see \autoref{fig:e02_f104_c16_layout}). 
		\begin{itemize}
			\item At T2\_Bnd\_B1\_x0m a constant discharge of 790.57 $m^3$/s is imposed.
			\item At T2\_Bnd\_B2\_x10000m a constant water level of 1.0 m is imposed.
			\item Branch B1 contains two "open Y-Z rectangular cross-sections" (width=\textbf{\textcolor[rgb]{0.5,0,0.15}{500}} m, Chezy=50 $m^{0.5}$/s), respectively having an elevation of \textbf{10.15} m AD at T2\_Bnd\_B1\_x0m and \textbf{10.00} m AD at T2\_ConNode. 
			\item Branch B2 contains two "open Y-Z rectangular cross-sections" (width=\textbf{\textcolor[rgb]{0.5,0,0.15}{500}} m, Chezy=50 $m^{0.5}$/s), respectively having an elevation \textbf{10.00} m at T2\_ConNode and \textbf{0.00} m AD at \newline T2\_Bnd\_B2\_x10000m.
		\end{itemize} 
		
		\Note The bathymetry and hydraulic properties of branch B1 and B2 of sub-model \textbf{\textcolor[rgb]{1,0,0}{T2}} are indentical to those branch B1 of sub-model \textbf{\textcolor[rgb]{1,0,0}{T1}}. \\
	
	\item Sub-model \textbf{\textcolor[rgb]{1,0,0}{T3}} comprises of two branches \textbf{B1} (length 15000 m) and \textbf{B2} (length 10000 m) , that are mutually connected at connection node T3\_ConNode (see \autoref{fig:e02_f104_c16_layout}).
		\begin{itemize}
			\item At T3\_Bnd\_B1\_x0m a constant discharge of 790.57 $m^3$/s is imposed.
			\item At T3\_Bnd\_B2\_x10000m a constant water level of 1.0 m is imposed.
			\item Branch B1 contains two "open Y-Z rectangular cross-sections" (width=\textbf{\textcolor[rgb]{0.5,0,0.15}{500}} m, Chezy=50 $m^{0.5}$/s), respectively having an elevation of \textbf{10.15} m AD at T3\_Bnd\_B1\_x0m and \textbf{10.00} m AD at T3\_ConNode. 
			\item Branch B2 contains two "open Y-Z rectangular cross-sections" (width=\textbf{\textcolor[rgb]{0.5,0,0.15}{500}} m, Chezy=50 $m^{0.5}$/s), respectively having an elevation of \textbf{10.00} m at T3\_ConNode and \textbf{0.00} m AD at \newline T3\_Bnd\_B2\_x10000m.
		\end{itemize} 
			
		\Note The bathymetry and hydraulic properties of branch B1 and B2 of sub-model \textbf{\textcolor[rgb]{1,0,0}{T3}} are indentical to branches B1 and B2 of sub-model \textbf{\textcolor[rgb]{1,0,0}{T2}}. The difference between \textbf{\textcolor[rgb]{1,0,0}{T2}} and \textbf{\textcolor[rgb]{1,0,0}{T3}} is the angle between the two branches connected at the connection node. \textbf{\textcolor[rgb]{1,0,0}{T2}} is a straight channel (angle of 180\textdegree).  \textbf{\textcolor[rgb]{1,0,0}{T3}} has, on the other hand, an angle of circa 143\textdegree \, between branch B1 and B2.\\
	
	\item Sub-model \textbf{\textcolor[rgb]{1,0,0}{T4}} comprises of four branches \textbf{B1} (length 15000 m), \textbf{B2} (length 10000 m), \textbf{B3} (length 10000 m), and \textbf{B4} (length 15000 m), that are mutually connected at connection node T4\_ConNode (see \autoref{fig:e02_f104_c16_layout}).
		\begin{itemize}
			\item At T4\_Bnd\_B1\_x0m a constant discharge of 790.57 $m^3$/s is imposed.
			\item At T4\_Bnd\_B2\_x10000m a constant water level of 1.0 m is imposed.
			\item Branch B1 contains two "open Y-Z rectangular cross-sections" (width=\textbf{\textcolor[rgb]{0.5,0,0.15}{500}} m, Chezy=50 $m^{0.5}$/s), respectively having an elevation of \textbf{10.15} m AD at T4\_Bnd\_B1\_x0m and \textbf{10.00} m AD at T4\_ConNode. 
			\item Branch B2 contains two "open Y-Z rectangular cross-sections" (width=\textbf{\textcolor[rgb]{0.5,0,0.15}{500}} m, Chezy=50 $m^{0.5}$/s), respectively having an elevation \textbf{10.00} m AD at T4\_ConNode and \textbf{0.00} m AD at \newline T4\_Bnd\_B2\_x10000m.
		\end{itemize}
		
			\Note The bathymetry and hydraulic properties of branch B1 and B2 of sub-model \textbf{\textcolor[rgb]{1,0,0}{T4}} are indentical to branches B1 and B2 of sub-model \textbf{\textcolor[rgb]{1,0,0}{T2}}. The difference between \textbf{\textcolor[rgb]{1,0,0}{T2}} and \textbf{\textcolor[rgb]{1,0,0}{T4}} is the angle between the two branches connected at the connection node. \textbf{\textcolor[rgb]{1,0,0}{T2}} is a straight channel (angle of 180\textdegree).  \textbf{\textcolor[rgb]{1,0,0}{T4}} has, on the other hand, an angle of circa 106\textdegree \, between branch B1 and B2.\\
	
	\end{enumerate}	
%---------------------------------------------------------------------------------
\section*{Results}
Water level, water depths, discharge and velocity results of sub-models \textbf{\textcolor[rgb]{1,0,0}{T1}} to \textbf{\textcolor[rgb]{1,0,0}{T4}} are given in \autoref{fig:e02_f104_c16_height} and \autoref{fig:e02_f104_c16_flow}.

\begin{figure}[H]
    \centering
    \includegraphics*[width=0.8\textwidth]{figures/e02_f104_c16_h1ands1.jpg}
    \caption{Water levels and water depths at (connection) nodes}
    \label{fig:e02_f104_c16_height}
\end{figure} 

\begin{figure}[H]
    \centering
    \includegraphics*[width=1.0\textwidth]{figures/e02_f104_c16_Q1andu1.jpg}
    \caption{Discharges and velocities at u-points directly adjacent to (connection) nodes}
    \label{fig:e02_f104_c16_flow}
\end{figure} 

%---------------------------------------------------------------------------------
\section*{Analysis of results}
The result in \autoref{fig:e02_f104_c16_height} and \autoref{fig:e02_f104_c16_flow} are analyzed by comparing water levels, water depths, discharges and velocities of \textbf{\textcolor[rgb]{1,0,0}{T3}} and \textbf{\textcolor[rgb]{1,0,0}{T4}} to \textbf{\textcolor[rgb]{1,0,0}{T2}}.

This comparisons shows the following :
\begin{enumerate}
\item No differences between \textbf{\textcolor[rgb]{1,0,0}{T1}} and \textbf{\textcolor[rgb]{1,0,0}{T2}} can be observed for these variables.
\item Discharges measured at u-points directly adjacent to (connection) nodes are approx the same, namely 790.57 $m^3$/s.
\item Water levels and water depths are higher for \textbf{\textcolor[rgb]{1,0,0}{T3}} and \textbf{\textcolor[rgb]{1,0,0}{T4}} than for \textbf{\textcolor[rgb]{1,0,0}{T2}}.  Next to this, a larger difference in water depths between sub-models \textbf{\textcolor[rgb]{1,0,0}{T2}} and \textbf{\textcolor[rgb]{1,0,0}{T4}} was found compared to the difference between  \textbf{\textcolor[rgb]{1,0,0}{T2}} and \textbf{\textcolor[rgb]{1,0,0}{T3}} (see \autoref{fig:e02_f104_c16_height}).  In other words, a smaller angle between branch B1 and B2 (sharper curved channel) results in an even higher water depth at the connection nodes. 
\item For both \textbf{\textcolor[rgb]{1,0,0}{T3}} and \textbf{\textcolor[rgb]{1,0,0}{T4}} the velocities became higher on the u-points directly in front of the connection node compared to \textbf{\textcolor[rgb]{1,0,0}{T2}} (see \autoref{fig:e02_f104_c16_flow}), whereas this was vice versa for the velocities at u-points directly behind the connection nodes. 
\end{enumerate}
%---------------------------------------------------------------------------------
\printrefsegment
